% Part: sets-functions-relations
% Chapter: sets
% Section: important-sets

\documentclass[../../../include/open-logic-section]{subfiles}

\begin{document}

\olfileid{sfr}{set}{imp}
\olsection{Sawetara Himpunan Wigati}

\begin{ex}
Umume kita bakal ngrembug himpunan kang !!{element}sne arupa
objek matematis. Ana papat himpunan kaya mangkono kang cukup
wigati nganti diwenehi jeneng khusus:
\begin{multline*}
    \Nat = \{0, 1, 2, 3, \ldots\} \\
    \shoveright{\text{himpunan wilangan natural}}\\
    \shoveleft{\Int = \{\ldots, -2, -1, 0, 1, 2, \ldots\}} \\
    \shoveright{\text{himpunan wilangan wutuh}}\\
    \shoveleft{\Rat = \Setabs{\nicefrac{m}{n}}{m, n \in \Int\text{ lan }n \neq 0}}\\
    \shoveright{\text{himpunan wilangan rasional}}\\
    \shoveleft{\Real = (-\infty, \infty)}\\
    \text{himpunan wilangan real (kontinuum)}
\end{multline*}
Kabeh iki himpunan \emph{tanpa wates}, yaiku saben himpunan
duwe !!{element}s kang cacahe tanpa wates.

Nalika pindhah saka himpunan siji menyang sabanjure ing dhaptar iki,
kita nambah \emph{luwih akeh} wilangan ing pasadhiyan kita.
Mesthine cetha yen $\Nat \subseteq \Int
\subseteq \Rat \subseteq \Real$: saben wilangan natural iku
wilangan wutuh; saben wilangan wutuh iku rasional; lan saben
wilangan rasional iku real. Semono uga, mesthine cetha yen
$\Nat \subsetneq \Int \subsetneq
\Rat$, amarga $-1$ iku wilangan wutuh nanging dudu wilangan natural,
lan $\nicefrac{1}{2}$ iku rasional nanging dudu wilangan wutuh.
Kang kurang cetha yaiku $\Rat \subsetneq \Real$, yaiku ana
wilangan real kang ora rasional\oliflabeldef{sfr:arith:real:realline}{,
nanging bab iki bakal dirembug maneh ing \olref[arith][real]{realline}}{}.

Kadhangkala kita uga nganggo himpunan wilangan wutuh positif
$\PosInt = \{1,2, 3, \dots\}$ lan himpunan kang mung ngemot
rong wilangan natural kapisan $\Bin = \{0, 1\}$.
\end{ex}

\begin{tagblock}{compsci}
\begin{ex}[Rerangken]
Tuladha liyane kang narik kawigaten yaiku himpunan $A^{*}$ saka
\emph{rerangken winates} ing alfabet $A$: saben urutan winates
anggota~$A$ iku rerangken ing $A$. Kita nglebokake
\emph{rerangken kosong $\Lambda$} ing antarane rerangken ing~$A$,
kanggo saben alfabet~$A$. Tuladhane,
\begin{multline*}
\Bin^*
=\{\Lambda,0,1,00,01,10,11,\\
000,001,010,011,100,101,110,111,0000,\ldots\}.
\end{multline*}
Yen $x=x_{1}\ldots x_{n}\in A^{*}$ iku rerangken kang dumadi
saka $n$ "aksara" saka $A$, mula kita nyebut \emph{dawane}
rerangken kasebut~$n$ lan nulis $\len{x}=n$.
\end{ex}
\end{tagblock}

\begin{ex}[Urutan tanpa wates]
Kanggo saben himpunan $A$, kita uga bisa ngrembug himpunan~$A^\omega$
saka urutan tanpa wates kang ngemot !!{element}s saka~$A$.
Urutan tanpa wates $a_1a_2a_3a_4\dots$ dumadi saka dhaptar
objek kang tanpa wates ing siji arah, lan saben objek iku
!!a{element} saka~$A$.
\end{ex}

\end{document}
